\chapter*{Conclusiones}
\addcontentsline{toc}{chapter}{Conclusiones}
\label{cap:conclusiones}

Se construyó un prototipo de sistema web que automatiza el análisis de los videos de la
prueba de nado forzado, desde la carga del video hasta los reportes, y se cumplió el
objetivo general. El sistema registra experimentos con sus grupos y tandas, recibe los
videos, los analiza en segundo plano con un \textit{pipeline} que localiza a cada espécimen
por sustracción de fondo y clasifica su conducta segundo a segundo con una red
convolucional espaciotemporal, permite revisar y corregir esa clasificación, y entrega
resultados y reportes en PDF, CSV y XLSX. El desempeño del clasificador se validó con
métricas de clasificación estándar, y el sistema cumple los atributos de calidad de la
norma ISO/IEC~25010 que se definieron como criterios de aceptación.

Se cumplieron los seis objetivos específicos. Se revisaron los trabajos y herramientas
existentes para el análisis automatizado de la prueba y se identificaron los enfoques de
visión por computadora y de aprendizaje supervisado más adecuados (OE-1). Se documentó el
protocolo experimental del laboratorio, con las conductas que se anotan, la forma de
grabar y las características de los videos disponibles (OE-2). Se diseñó la arquitectura
conforme a la norma ISO/IEC/IEEE~12207, con los módulos de carga de videos,
\textit{pipeline} de análisis, base de datos y visualización de resultados (OE-3). Se
implementó el módulo de análisis conductual, con la localización de cada espécimen y un
clasificador entrenado con segundos etiquetados por el equipo (OE-4). Se evaluó el
desempeño del clasificador sobre un conjunto de prueba independiente, con precisión,
\textit{recall} y F1 por conducta, y se alcanzó la meta de F1 por conducta de RNF-02
(OE-5). Y se cuantificó la reducción del tiempo de análisis frente al proceso manual,
además de dejar documentado el diseño modular del módulo de clasificación y qué tendría
que cambiar el \textit{pipeline} para aplicarlo a otro paradigma conductual (OE-6).

El alcance del trabajo tiene límites que conviene tener presentes. El clasificador se
entrenó y se validó solo para la prueba de nado forzado, con videos de un solo
laboratorio, y el sistema asume una calidad mínima de video: no corrige reflejos, pelaje
claro ni contraste insuficiente, y cuando no encuentra los cilindros se detiene y genera
un reporte de diagnóstico. El video del día 1, que es una sesión de habituación, se
guarda como constancia y no se analiza. Estos límites no restan al cumplimiento de los
objetivos: son las fronteras que el alcance fijó desde el inicio.

En conclusión, el sistema es más que una herramienta de software: es un apoyo concreto
para la investigación del Laboratorio de Bioquímica Estructural y para la institución que
lo alberga. Al centralizar los videos, los resultados y los reportes de cada experimento,
deja la evidencia conservada y consultable, y libera al personal de horas de anotación
manual para dedicarlas al diseño y la interpretación de los experimentos. Su mayor valor
está en la búsqueda de compuestos con posible efecto antidepresivo: la prueba de nado
forzado es uno de los modelos preclínicos con que se comparan los compuestos candidatos
contra un antidepresivo de referencia como la fluoxetina, y un análisis más rápido,
reproducible y revisable por el investigador permite evaluar más compuestos y más dosis
en el mismo tiempo, con resultados comparables entre grupos y entre sesiones. El sistema
no diagnostica ni sustituye el criterio del investigador: entrega una clasificación por
segundo que el laboratorio revisa, corrige y usa para decidir qué compuestos merecen
avanzar en el desarrollo de nuevos tratamientos contra la depresión.
